% Drawing 29/10-1 -- folder 29, batch 10, archive page 183 (PDF page 184).
% Transcription: en haut à gauche de la feuille, un arbre orienté : d'une racine à gauche part une arête vers un premier sommet, qui se ramifie vers la droite en deux branches, elles-mêmes ramifiées ; toutes les arêtes sont fléchées en s'éloignant de la racine, les sommets marqués de points.
% Kind: rooted tree with arrows away from the root. Method 2 (schematic redraw from the facsimile). Status: ok.
% Notes: As on the page: the root at the left, one edge to the first vertex, which branches into an upper and a lower branch; the upper branch divides again (a heavy branch going up to the right, which forks at its end, and a branch to the right, which forks twice); the lower branch runs to the right through two vertices and forks at the second. Eighteen edges, each drawn as an arrow between its two vertices, pointing away from the root; the vertices are marked with points (leaves are arrow tips). This figure is not mentioned in the notes of the transcription. Overlaid 2026-11-12 (Opus 5.5) on the 400 dpi colour crop: the eleven vertices and eighteen edges fall on the page; corrected: on the edges that end at a vertex the arrowhead stands before the vertex, as on the page (about four fifths along; on the root edge half-way), not at it; leaf edges end in their arrowhead. Redrawn 2026-11-12 (issue 128): to check again.
% Coordinates: 1 unit = 1 mm = 1/100 of the crop width (bbox in page fractions: [0.11, 0.155, 0.34, 0.28]).
% Redrawn by Claude (claude-fable-5-2), issue #128, 2026-11-12; not checked by a human; not compiled (no TeX in the session).
% Overlaid by Claude (claude-opus-5-5) 2026-11-12 (issue #128): rasterised from its own coordinates over the 400 dpi colour facsimile and corrected from it. Not compiled. Status: ok.
% Rights: reproduces part of the work -- behind RENDER_TRANSCRIPTS, never published by itself.
\documentclass[tikz,border=3pt]{standalone}
\usepackage{amsmath,amssymb}
\usetikzlibrary{patterns,arrows.meta}
\tikzset{pencil/.style={black!45}, struck/.style={black!35}, mpt/.style={circle,draw,inner sep=1.3pt}, pt/.style={circle,fill,inner sep=0.9pt}, opt/.style={circle,draw,fill=white,inner sep=1.1pt}}
\begin{document}
\begin{tikzpicture}[x=1mm,y=1mm,line cap=round,line join=round,font=\scriptsize]
\draw (9,-47) -- (33,-47.5);
\draw[-{Stealth[length=2mm]}] (9,-47) -- (22.2,-47.3);
\draw (33,-47.5) -- (44.5,-40.5);
\draw[-{Stealth[length=2mm]}] (33,-47.5) -- (42.2,-41.9);
\draw (33,-47.5) -- (43,-54.5);
\draw[-{Stealth[length=2mm]}] (33,-47.5) -- (41.0,-53.1);
\draw[very thick] (44.5,-40.5) -- (55,-30.5);
\draw[-{Stealth[length=2mm]},very thick] (44.5,-40.5) -- (52.9,-32.5);
\draw[very thick] (55,-30.5) -- (66,-20.5);
\draw[-{Stealth[length=2mm]},very thick] (55,-30.5) -- (63.8,-22.5);
\draw[-{Stealth[length=2mm]}] (66,-20.5) -- (76,-8);
\draw[-{Stealth[length=2mm]}] (66,-20.5) -- (82,-16);
\draw (44.5,-40.5) -- (58,-41);
\draw[-{Stealth[length=2mm]}] (44.5,-40.5) -- (55.3,-40.9);
\draw (58,-41) -- (68,-32);
\draw[-{Stealth[length=2mm]}] (58,-41) -- (66.0,-33.8);
\draw[-{Stealth[length=2mm]}] (68,-32) -- (78,-23);
\draw[-{Stealth[length=2mm]}] (68,-32) -- (81,-32.5);
\draw (58,-41) -- (74,-41.5);
\draw[-{Stealth[length=2mm]}] (58,-41) -- (70.8,-41.4);
\draw[-{Stealth[length=2mm]}] (74,-41.5) -- (87.5,-36);
\draw[-{Stealth[length=2mm]}] (74,-41.5) -- (89,-44);
\draw (43,-54.5) -- (59,-55.5);
\draw[-{Stealth[length=2mm]}] (43,-54.5) -- (55.8,-55.3);
\draw (59,-55.5) -- (70.5,-56);
\draw[-{Stealth[length=2mm]}] (59,-55.5) -- (68.2,-55.9);
\draw[-{Stealth[length=2mm]}] (70.5,-56) -- (83,-56);
\draw[-{Stealth[length=2mm]}] (70.5,-56) -- (84,-68.5);
\node[pt] at (9,-47) {};
\node[pt] at (33,-47.5) {};
\node[pt] at (44.5,-40.5) {};
\node[pt] at (43,-54.5) {};
\node[pt] at (55,-30.5) {};
\node[pt] at (66,-20.5) {};
\node[pt] at (58,-41) {};
\node[pt] at (68,-32) {};
\node[pt] at (74,-41.5) {};
\node[pt] at (59,-55.5) {};
\node[pt] at (70.5,-56) {};
\end{tikzpicture}
\end{document}
